\clearpage
\section*{Anexo 4.1-S --- Puntos de vista y correspondencias}
\phantomsection
\label{anx:S}
\addcontentsline{toc}{section}{Anexo 4.1-S --- Puntos de vista y correspondencias}
La entidad de interés es la plataforma logística Puelche y sus fronteras con ERP, personas y servicios de terceros. La descripción cubre toma de pedido, reserva, preparación, salida, entrega, rendición, trazabilidad y análisis; su horizonte incluye dos etapas y 36 meses de operación. Cada punto de vista fija preocupación, interesados, modelos y comprobación, conforme a RT-02.03 e ISO/IEC/IEEE 42010:2022 (ISO, 2022a).

\begin{tablalafrox}{Interesados y modelos de las cinco vistas}{tab:puntos-vista}{F{0.15} F{0.19} F{0.29} Y}{\cab{Vista} & \cab{Interesados} & \cab{Preocupación} & \cab{Modelos y evidencia}}
Lógica & TI y responsables de módulo. & Dueños, límites e interfaces sin dependencia de tablas ajenas. & General, capas y catálogos D/F/G/H/N. \\
Procesos & Operaciones, conductores y Tesorería. & Una captura avanza a confirmación, salida y rendición con estados claros. & Secuencias ARQL-15/16; eventos A y documento U. \\
Despliegue & TI, Infraestructura y Operaciones. & Cada función dispone de realización y contingencia por sitio. & Inventario N y enlaces exactos a 4.2--4.3. \\
Datos & Calidad, Datos y Comercial. & Lote, reservas, acuses y evidencias tienen dueño y retención. & Modelo ARQL-18, reconciliación K y capítulo 5. \\
Seguridad & Seguridad, Calidad y CLIENTE. & Nadie libera, cobra o consulta sin permiso y evidencia. & Límites de confianza, Q/R y prueba AL-OFF-01. \\
\end{tablalafrox}
{\footnotesize Fuente: interesados del caso, RT-02.03 y organización de la descripción arquitectónica (ISO, 2022a).\par}

CR-01 exige que cada capacidad del esquema y explicación de solución tenga uno o más componentes N con nombre idéntico y que todo componente tenga capacidad de origen. CR-02 exige productor único, consumidores registrados y contrato por evento; una arista sin contrato se rechaza. CR-03 exige que cada escritura tenga dueño de datos, clave idempotente y auditoría. CR-04 exige para cada función desconectada persistencia, permisos, límite de autonomía y reconciliación. CR-05 vincula componente lógico a realización física exacta; no se cuenta una referencia genérica. CR-06 exige que cada amenaza tenga control y prueba, y que cada decisión ADR tenga consecuencia verificable.

Arquitectura mantiene un registro de diferencias con componente, regla infringida, requisito, dueño, decisión y evidencia de cierre. La comprobación se ejecuta al cambiar un módulo, contrato, tecnología o etapa. AL-MAP-01 verifica CR-01/05; pruebas de contrato, CR-02/03; AL-OFF-01, CR-04; y la revisión de Q/R/O, CR-06. La conformidad de una vista no se atribuye a otra sin sus modelos.
